% Auto-ABM --- Related Work (draft)
% Format follows the EMNLP/ACL template in ../emnlp2026 (acl.sty + acl_natbib.bst).
% Compile:  pdflatex related_work && bibtex related_work && pdflatex related_work && pdflatex related_work

\documentclass[11pt]{article}

\usepackage[review]{acl}
\usepackage{times}
\usepackage{latexsym}
\usepackage[T1]{fontenc}
\usepackage[utf8]{inputenc}
\usepackage{microtype}
\usepackage{inconsolata}

\title{Auto-ABM: Related Work (Draft)}

\author{Anonymous \\
  \texttt{anonymous@example.com}}

\begin{document}
\maketitle

\section{Related Work}

\subsection{The Agent-Based Modeling Ecosystem}

Agent-based modeling (ABM) represents a system as autonomous agents that follow
local rules and interact with one another and their environment. The modeler
specifies individual behavior and observes how macro patterns emerge from the
bottom up \citep{bonabeau2002abm,macal2010tutorial}. Classic studies such as
Schelling's segregation model \citep{schelling1971segregation} and Epstein and
Axtell's Sugarscape \citep{epstein1996growing} established ABM as a standard
tool across ecology, epidemiology, economics, and the social sciences. A mature
software ecosystem supports this work: NetLogo \citep{wilensky1999netlogo},
Repast Simphony \citep{north2013repast}, MASON \citep{luke2005mason}, Mesa
\citep{kazil2020mesa,terhoeven2025mesa3}, and GAMA \citep{taillandier2019gama}.
The ODD protocol provides a common standard for documenting models
\citep{grimm2006odd,grimm2020odd}.

Building and studying a model nevertheless remains largely manual work that
requires both programming skill and domain knowledge. Researchers still move by
hand through reading literature, writing an ODD, translating it into code,
debugging, sweeping parameters, and interpreting output
\citep{macal2010tutorial,grimm2020odd}. Recent work has brought large language
models (LLMs) into isolated steps of this cycle. NetLogo Chat assists with
writing and learning NetLogo code \citep{chen2024netlogochat}; few-shot prompting
and retrieval improve generated code accuracy \citep{martinez2024netlogo}; and
SAGE uses a verifier to iteratively produce runnable models
\citep{niu2024sage}. Connectors built on the Model Context Protocol
\citep{mcp2024}, such as NetLogo-MCP \citep{netlogomcp2025}, let an assistant
create and run models in conversation. Each tool automates a single chore rather
than the wider loop in which the simulation, its execution trace, and its ODD
are the objects of study.

\subsection{AI-Driven Automation of the Modeling Workflow}

A separate line of research asks how far LLM agents can automate multi-step
scientific work. ReAct interleaves reasoning with tool use in an iterative loop
\citep{yao2023react}; wrapping an LLM with tools, memory, and control logic in
an agentic harness yields systems that can plan and execute extended workflows
\citep{wang2024survey}. This pattern now drives broader AutoResearch efforts.
The AI Scientist automates ideation, coding, experimentation, and writing for
machine-learning research \citep{lu2024aiscientist}; Coscientist plans and runs
real chemistry experiments through laboratory automation
\citep{boiko2023coscientist}; and AgentEconomist translates economic intuitions
into executable experiments over a fixed simulator, grounding each stage in the
literature \citep{agenteconomist2026}. These systems span a full study, but they
target program code, wet-lab reactions, or one domain-specific simulator rather
than the simulation, trace, and ODD that define an ABM research workflow.

Closer to our setting, several systems automate parts of the modeling pipeline
from natural language. SAGE and SOCIA generate and repair simulator code from
scenario descriptions \citep{niu2024sage,hua2025socia}; AgentEconomist
orchestrates hypothesis formation, experiment design, and execution over a
simulator \citep{agenteconomist2026}. Auto-ABM pursues a broader goal: rather
than producing one-off code or automating a single stage, it orchestrates the
full research loop around a classical, deterministic ABM. To our knowledge, no
prior system combines natural-language workflow orchestration, deterministic
execution on a reusable kernel, ODD-aligned documentation, trace-grounded
explanation, and one-click reproducible export in a single simulation-first
workbench.

\bibliography{custom}

\end{document}
